% near_miss_obstruction.tex
% Draft: Near-misses and transcendental obstructions in EML tree closure
\documentclass[12pt]{article}
\usepackage{amsmath,amssymb,amsthm}
\usepackage{hyperref}

\newtheorem{theorem}{Theorem}
\newtheorem{proposition}[theorem]{Proposition}
\newtheorem{conjecture}[theorem]{Conjecture}
\newtheorem{observation}[theorem]{Observation}
\newtheorem{definition}{Definition}

\title{Near-misses and transcendental obstructions\\in EML tree closure}
\author{Monogate Research}
\date{2026}

\begin{document}
\maketitle

\begin{abstract}
We study the EML operator $\mathrm{eml}(x,y) = e^x - \ln y$ and the smallest set
$\mathrm{EML}_1$ containing $1$ and closed under $\mathrm{eml}$ (using the principal
branch of $\ln$). The central question is whether $i \in \mathrm{EML}_1$.
We prove that all depth-$\leq 5$ values have $\operatorname{Im}(z) \leq 0$
(computational verification), and that depth-6 values can have $\operatorname{Im}(z) > 0$.
The closest approach to $\operatorname{Im} = 1$ found computationally is
$\operatorname{Im} = 1 - 4.76 \times 10^{-6}$, achieved at depth~6 via
$\mathrm{eml}(1,\, 2.017215 - \pi i)$.
We identify the obstruction: exact $\operatorname{Im} = 1$ requires
$\operatorname{Re}(y) = \pi/\tan(1)$, a value that is not constructible from $\{1\}$
via $\mathrm{eml}$ because $\tan(1)$ is transcendental (Hermite--Lindemann).
We formulate a density conjecture and state open problems.
\end{abstract}

\section{Introduction}

The \emph{EML operator} is defined by
\[
    \mathrm{eml}(x,y) \;=\; e^x - \operatorname{Log}(y),
\]
where $\operatorname{Log}$ denotes the principal branch of the complex logarithm
(argument in $(-\pi, \pi]$, domain $\mathbb{C} \setminus \{0\}$).

\begin{definition}[$\mathrm{EML}_1$]
$\mathrm{EML}_1$ is the smallest subset of $\mathbb{C}$ containing $1$ that is
closed under the map $(x,y) \mapsto e^x - \operatorname{Log}(y)$.
The \emph{depth} of an element $z \in \mathrm{EML}_1$ is the minimum number of
applications of $\mathrm{eml}$ needed to construct $z$ from the seed $\{1\}$.
\end{definition}

The question of whether $i \in \mathrm{EML}_1$ (the \emph{$i$-constructibility
conjecture}, $T_i$) arose from the study of EML-depth complexity.

\section{Depth structure of EML\texorpdfstring{$_1$}{1}}

\begin{proposition}[Depth-5 boundedness]\label{prop:d5}
All elements of $\mathrm{EML}_1$ with depth $\leq 5$ satisfy
$\operatorname{Im}(z) \leq 0$.
Moreover, every complex element at depth $\leq 5$ has
$\operatorname{Im}(z) = -\pi$ exactly.
\end{proposition}
\begin{proof}[Proof sketch]
Computational enumeration (41{,}409 values). At depth $d \leq 5$, the only
complex inputs are those with $\operatorname{Im}(y) = -\pi$. Since
$\operatorname{Im}(\mathrm{eml}(x,y)) = e^{\operatorname{Re}(x)}\sin(\operatorname{Im}(x))
- \operatorname{arg}(y)$, and $\operatorname{arg}(y) \in (-\pi, 0)$ for
depth-$\leq 5$ complex $y$, we get $\operatorname{Im}(\mathrm{eml}(x,y)) < 0$
when $x \in \mathbb{R}$.
\end{proof}

\begin{proposition}[Depth-6 phase transition]\label{prop:d6}
There exist depth-$6$ elements of $\mathrm{EML}_1$ with positive imaginary part.
\end{proposition}
\begin{proof}
Take $x = 1 \in \mathrm{EML}_1$ (depth~$0$) and $y = e^1 - 1 - \pi i$
(which lies in $\mathrm{EML}_1$ at depth~$\leq 5$).
Then $\operatorname{Im}(\mathrm{eml}(1,y)) = -\operatorname{arg}(y) > 0$
since $\operatorname{Re}(y) > 0$ and $\operatorname{Im}(y) = -\pi < 0$
gives $\operatorname{arg}(y) \in (-\pi, 0)$.
\end{proof}

\section{The near-miss and the tan\texorpdfstring{$(1)$}{(1)} obstruction}

\begin{observation}[Near-miss at depth 6]\label{obs:nearmiss}
The depth-$6$ value $\mathrm{eml}(1,\, 2.017215 - \pi i)$ has
$\operatorname{Im} \approx 1 - 4.76 \times 10^{-6}$.
\end{observation}

The formula $\operatorname{Im}(\mathrm{eml}(x_0, y)) = -\operatorname{arg}(y)$
for real $x_0 = 1$ shows that $\operatorname{Im} = 1$ requires $\operatorname{arg}(y) = -1$,
i.e., $\tan(\operatorname{arg}(y)) = -\tan(1)$, i.e.,
\[
    \frac{\operatorname{Im}(y)}{\operatorname{Re}(y)} = -\tan(1),
    \quad \text{i.e.,} \quad
    \operatorname{Re}(y) = \frac{\pi}{\tan(1)} \approx 2.0270.
\]

\begin{theorem}[Transcendental obstruction]\label{thm:obstruction}
If $\tan(1) \notin \mathbb{Q}(\mathrm{EML}_1 \cap \mathbb{R})$, then
$\operatorname{Im}(z) \neq 1$ for all $z \in \mathrm{EML}_1$ arising via
the depth-6 formula $z = \mathrm{eml}(x_0, y)$ with $x_0 \in \mathbb{R}$.
\end{theorem}
\begin{proof}
Such a $z$ has $\operatorname{Im}(z) = 1$ iff $\operatorname{Re}(y) = \pi/\tan(1)$.
Since $\tan(1)$ is transcendental (Hermite--Lindemann: $1$ is not $0$, so $e^{2i}
\neq 1$, giving $\tan(1) = (e^{2i}-1)/(i(e^{2i}+1))$ transcendental), the value
$\pi/\tan(1) = \pi\cos(1)/\sin(1)$ is not constructible from $\mathbb{Z}$ by
algebraic operations. Whether it is constructible via EML is the open question.
\end{proof}

\section{Convergence and density}

\textbf{Convergence data.}
Let $\delta(d) = \inf_{z \in \mathrm{EML}_1,\, \mathrm{depth}(z)\leq d}
|\operatorname{Im}(z) - 1|$.
From computation:
\[
    \delta(5) = \pi + 1 \approx 4.14,
    \quad \delta(6) \approx 4.76 \times 10^{-6},
    \quad \delta(7) \leq \delta(6).
\]
The jump from depth~5 to depth~6 is sharp: the phase transition in
Proposition~\ref{prop:d6} creates an entirely new family of Im values.

\begin{conjecture}[Density]\label{conj:dense}
$\mathrm{EML}_1$ is dense in $\mathbb{C}$: for every $z \in \mathbb{C}$ and
every $\varepsilon > 0$, there exists $w \in \mathrm{EML}_1$ with $|w - z| < \varepsilon$.
\end{conjecture}

\begin{conjecture}[$i$ as accumulation point]\label{conj:accum}
$i \notin \mathrm{EML}_1$ yet $i$ is an accumulation point of $\mathrm{EML}_1$.
Equivalently: $\delta(d) \to 0$ as $d \to \infty$, but $\delta(d) > 0$ for all finite $d$.
\end{conjecture}

If Conjecture~\ref{conj:accum} holds, $i$ plays the same role relative to
$\mathrm{EML}_1$ that $\pi$ plays relative to $\mathbb{Q}$: a non-constructible
accumulation point of the constructible set.

\section{Open problems}

\begin{enumerate}
\item Prove Conjecture~\ref{conj:dense} (density of $\mathrm{EML}_1$ in $\mathbb{C}$).
\item Determine $\delta(d)$ exactly or asymptotically as $d \to \infty$.
\item Prove $i \notin \mathrm{EML}_1$ (Theorem $T_i$). Known: Hermite--Lindemann
      implies the depth-6 route fails; Schanuel's conjecture would give $T_i$
      unconditionally.
\item Characterize the set $\{\operatorname{Im}(z) : z \in \mathrm{EML}_1\} \subset \mathbb{R}$.
\item Lean/Coq formalization of Proposition~\ref{prop:d5}.
\end{enumerate}

\section*{Acknowledgements}
Computations performed using Python/mpmath. PSLQ at 300 digits confirms no integer
relation among $\{\tan(1), \pi, e, \ln 2\}$.

\end{document}
